\subsection{Арифметические операции}
Операции: \verb|+ - * / % ^| (плюс, минус, умножение, деление, остаток, возведение в степень). Сравнения: \verb|< <= > >= == ~=|. Оператора \verb|**| в Lua нет ни в одной версии языка — это оператор из Python; для возведения в степень всегда используйте \verb|a ^ b|. Оператор целочисленного деления \verb|//| в симуляторе (Lua 5.3) доступен и работает, однако на других сборках Lua он может отсутствовать, поэтому для переносимости кода надёжнее писать \verb|math.floor(a/b)|. В Lua 5.3 у \hyperlink{term:number}{чисел} есть два подвида — целые (\verb|integer|) и вещественные (\verb|float|), — но для задач этого курса это различие почти не важно: арифметика работает с ними единообразно. Единственное, что стоит запомнить: оператор \verb|/| всегда возвращает вещественное число, даже если оба операнда целые (\verb|10 / 2| даёт \verb|5.0|, а не \verb|5|); поэтому для округлений применяются \verb|math.floor|, \verb|math.ceil| и \verb|math.abs|.

Приоритет: сначала выполняются \verb|^|, затем \verb|* / %|, далее \verb|+ -|. Скобки меняют порядок вычислений. Остаток \verb|%| определён как «остаток от деления», полезен для задач чётности и циклических индексов.
\begin{codefile}{lua/examples/04_language/lua_arithmetic_ops.lua}
\end{codefile}
\paragraph{Пояснения.}
- \verb|math.sqrt| возвращает корень квадратный; \verb|math.pi| — число $\pi$.
- \verb|math.max| и \verb|math.min| ограничивают значения: \verb|v = math.max(0, math.min(1, v))|.
- Деление на ноль следует избегать: проверяйте знаменатель перед делением.
- Перевод градусов в радианы и обратно можно писать вручную (\verb|deg * math.pi / 180|) или короче — через готовые функции \verb|math.rad(deg)| и \verb|math.deg(rad)|. Дальше в курсе (полёт по окружности, курс \verb|yaw|) встречаются оба варианта записи.
\subsection*{Задачи для самостоятельного решения}
\begin{enumerate}
  \item Вычислите выражение \verb|3 + 4 * 2 - 5|.
  \item Найдите остаток \verb|17 % 4|.
  \item Запишите \verb|3 ^ 4|.
  \item Вычислите среднее \verb|(x + y) / 2|.
  \item Запишите \verb|math.abs(x2 - x1)|.
  \item Вычислите \verb|math.floor(n / 3)|.
  \item Запишите условие «\verb|n| чётное» с \verb|%|.
  \item Вычислите \verb|math.sqrt(x*x + y*y)| (евклидова длина).
  \item Ограничьте \verb|v| в диапазоне \verb|[0,1]|: \verb|v = math.max(0, math.min(1, v))|.
  \item Переведите градусы в радианы: \verb|rad = deg * math.pi / 180|.
  \item Вычислите площадь прямоугольника \verb|S = w * h|.
  \item Запишите среднее трёх чисел \verb|(a+b+c)/3|.
  \item Реализуйте линейную интерполяцию \verb|lerp = a + t*(b-a)|.
  \item Посчитайте сумму массива \verb|a|: примените цикл \verb|for|.
  \item Найдите максимум массива \verb|a|: пройдитесь и сравнивайте.
\end{enumerate}
